problem:  
Let \( M^5 \) be a closed, simply connected Riemannian 5‑manifold admitting metrics of **nonnegative sectional curvature**.  
For any such metric \( g \), denote its isometry group by \( \mathrm{Iso}(M,g) \), and define the **symmetry rank**  
\[
\rho(g) := \mathrm{rank}\!\big(\mathrm{Iso}(M,g)\big).
\]  
It is known that \( \rho(g)\le3 \) for all \(g\).  A metric realizing \( \rho(g)=3 \) is called **maximally symmetric** under nonnegative curvature, while \( \rho(g)=2 \) is *almost maximal*.  

**Open problem:**  
Does there exist a simply connected 5‑manifold \(M\) admitting both a nonnegatively curved metric \( g_0 \) with symmetry rank \(2\) and another \( g_1 \) with rank \(3\), such that **no smooth one‑parameter family of nonnegatively curved metrics** \( g_t\ (t\in[0,1]) \) on \(M\) exists with the property that  
\[
g_t\to g_1 \text{ as } t\to 1,\quad \text{and} \quad \mathrm{Iso}(M,g_t)
\text{ contains (a subgroup conjugate to) }T^2 \ \forall t<1?
\]  
In words: can the passage from almost maximal to maximal symmetry rank fail to occur through any continuous deformation within the space of nonnegatively curved metrics on \(M\)?  
Equivalently, are there 5‑manifolds for which **symmetry enhancement under nonnegative curvature is topologically disconnected** in the moduli of metrics?

Why is it a "good" problem:  
- **Profound insight:** The problem asks whether **symmetry in nonnegatively curved geometry can change continuously**, or whether curvature imposes rigidity barriers between discrete symmetry levels.  This tests the interplay between metric deformation, curvature, and transformation group structure—the boundary of current understanding between geometry and Lie group representation on manifolds.  
- **Bridge between fields:** It connects **Riemannian deformation theory**, **transformation group rigidity**, and the **topology of moduli spaces** of nonnegatively curved metrics.  Resolving it would require synthesizing ideas from Alexandrov geometry (equivariant limits), Lie group actions, and smooth moduli theory, potentially extending Cheeger–Fukaya collapse theory to a symmetry‑sensitive setting.  
- **Extreme simplicity:**  Its statement is short and intuitive:  
  “Can symmetry under nonnegative curvature jump discretely between ranks 2 and 3, without a continuous deformation path?”  
  Yet this question encapsulates deep structural uncertainty about how curvature and symmetry interact in dimension five, making it both accessible and profoundly challenging.